\section{\IfLanguageName{dutch}{API laag}{API layer}}
\label{ch:architecture-api}

De API maakt deel uit van de solution en is ontwikkeld als een ASP.NET Core Razor Pages applicatie. Binnen deze applicatie worden specifieke API-endpoints voorzien voor data-access, waarbij uitsluitend GET-requests worden ondersteund. Hierdoor wordt gegarandeerd dat de API enkel gebruikt wordt voor het ophalen van data en niet voor het aanpassen ervan.

De API kan benaderd worden via HTTP-requests en is tevens beschikbaar binnen het afgesloten Docker-netwerk, zodat interne services zoals Jupyter Notebook-containers er veilig gebruik van kunnen maken.

Authenticatie wordt voorzien binnen de Razor Pages applicatie, bijvoorbeeld via cookies of OpenID Connect (OIDC). Dit zorgt ervoor dat enkel geautoriseerde gebruikers en services toegang krijgen tot de beschikbare endpoints.

De API fungeert als een gecontroleerde tussenlaag tussen de database en de eindgebruikers, waarbij directe toegang tot de databronnen vermeden wordt. De API communiceert met de database via een read-only gebruiker of via een read-replica, afhankelijk van de architectuur en performantievereisten.

Binnen de API worden verschillende maatregelen genomen om veiligheid, performantie en betrouwbaarheid te garanderen. Queryparameters worden uitgebreid gevalideerd om foutieve of kwaadaardige requests te vermijden. Daarnaast worden filters toegepast om enkel relevante data op te halen en wordt paginatie gebruikt om de hoeveelheid teruggegeven data per request te beperken.

Elke request wordt gelogd naar logbestanden, zodat monitoring, debugging en auditing mogelijk zijn.

Indien nodig kunnen gevoelige gegevens (zoals persoonsgegevens) bijkomend gefilterd of gemaskeerd worden (bijvoorbeeld via PII masking), zodat deze niet zichtbaar zijn voor eindgebruikers of analysetoepassingen.
